\documentclass[11pt,a4paper]{article}
% =====================================================================
%  Shared preamble for the Part V lecture notes (lectures 19-22).
%
%  These four lectures sit outside Geron, so the notes ARE the primary
%  source and are examined on the same terms as the textbook parts.
%  Everything here is chosen to make that readable on paper: one column,
%  generous measure, and the same six-colour palette as the slides so a
%  student moving between the two is not re-learning what red means.
%
%  Build with pdflatex (twice, for the toc and the refs):
%      cd notes && latexmk -pdf lecture-19.tex
%  or  make -C notes
% =====================================================================
\usepackage[T1]{fontenc}
\usepackage[utf8]{inputenc}
\usepackage{newtxtext,newtxmath}      % Times-like: prints well, reads well
\usepackage[scaled=0.86]{inconsolata} % code, at a size that sits beside Times
\usepackage{microtype}
% newtx defines \Bbbk, \openbox, \proof and \endproof; amssymb and amsthm
% then redefine them and abort. Freeing the four names before those two
% packages load is the standard fix, and it costs nothing here: we use
% amsthm only for \newtheorem.
\let\Bbbk\relax
\usepackage{amsmath,amssymb}
\let\openbox\relax
\let\proof\relax
\let\endproof\relax
\usepackage{amsthm}
\usepackage{booktabs}
\usepackage{enumitem}
\usepackage{xcolor}
\usepackage{tcolorbox}
\usepackage{titlesec}
\usepackage{graphicx}
\usepackage[hidelinks]{hyperref}
\tcbuselibrary{skins,breakable}

% ---- the course palette, from assets/css/custom.css -------------------
\definecolor{aimlprimary}{HTML}{0B3D62}   % deep blue  - structure
\definecolor{aimlaccent} {HTML}{C0392B}   % brick red  - failures
\definecolor{aimlsuccess}{HTML}{14663A}   % green      - repairs
\definecolor{aimlmath}   {HTML}{6C3483}   % purple     - the derivation
\definecolor{aimlmuted}  {HTML}{4B5563}
\definecolor{aimlrule}   {HTML}{B0BCC7}
\definecolor{aimlsoft}   {HTML}{F4F7F9}

\usepackage[a4paper,margin=32mm,top=30mm,bottom=30mm]{geometry}
\setlength{\parskip}{0.45\baselineskip}
\setlength{\parindent}{0pt}
\linespread{1.06}

% ---- headings --------------------------------------------------------
\titleformat{\section}{\Large\bfseries\color{aimlprimary}}{\thesection}{0.7em}{}
\titleformat{\subsection}{\large\bfseries\color{aimlprimary}}{\thesubsection}{0.7em}{}
\titleformat{\subsubsection}{\bfseries\color{aimlmuted}}{}{0em}{}
\titlespacing*{\section}{0pt}{1.6\baselineskip}{0.5\baselineskip}
\titlespacing*{\subsection}{0pt}{1.2\baselineskip}{0.35\baselineskip}

% ---- boxes -----------------------------------------------------------
% One box per role, and the role is the colour. A student who has seen the
% slides already knows what each one means before reading a word of it.
\newtcolorbox{keybox}[1][]{
  breakable, enhanced, colback=aimlsoft, colframe=aimlprimary,
  boxrule=0.4pt, left=8pt, right=8pt, top=6pt, bottom=6pt,
  fonttitle=\bfseries\small, coltitle=white, colbacktitle=aimlprimary,
  attach boxed title to top left={yshift=-2mm, xshift=4mm},
  boxed title style={boxrule=0pt, arc=1pt}, #1}

\newtcolorbox{failbox}[1][]{
  breakable, enhanced, colback=white, colframe=aimlaccent,
  boxrule=0.9pt, left=8pt, right=8pt, top=6pt, bottom=6pt,
  fonttitle=\bfseries\small, coltitle=white, colbacktitle=aimlaccent,
  attach boxed title to top left={yshift=-2mm, xshift=4mm},
  boxed title style={boxrule=0pt, arc=1pt}, #1}

\newtcolorbox{derivbox}[1][]{
  breakable, enhanced, colback=white, colframe=aimlmath,
  boxrule=0.6pt, left=8pt, right=8pt, top=6pt, bottom=6pt,
  fonttitle=\bfseries\small, coltitle=white, colbacktitle=aimlmath,
  attach boxed title to top left={yshift=-2mm, xshift=4mm},
  boxed title style={boxrule=0pt, arc=1pt}, #1}

% An exercise. Deliberately the same shape as the notebook's `try` field:
% one change, and what should happen to the output.
\newtcolorbox{trybox}[1][]{
  breakable, enhanced, colback=white, colframe=aimlsuccess,
  boxrule=0.5pt, left=8pt, right=8pt, top=6pt, bottom=6pt,
  fonttitle=\bfseries\small, coltitle=white, colbacktitle=aimlsuccess,
  attach boxed title to top left={yshift=-2mm, xshift=4mm},
  boxed title style={boxrule=0pt, arc=1pt}, #1}

% ---- small semantics -------------------------------------------------
\newcommand{\scope}[1]{\textcolor{aimlmuted}{\small #1}}
\newcommand{\term}[1]{\textbf{#1}}
\newcommand{\code}[1]{\texttt{\small #1}}
\newcommand{\lecture}[1]{Lecture~#1}

\newtheorem{definition}{Definition}[section]
\newtheorem{proposition}[definition]{Proposition}

% ---- title block -----------------------------------------------------
% \notestitle{19}{Information retrieval: the lexical foundation}{SciFact (BEIR)}
\newcommand{\notestitle}[3]{%
  \thispagestyle{empty}
  \begin{flushleft}
    {\color{aimlmuted}\small\sffamily
     APPLICATIONS OF MACHINE LEARNING \quad$\cdot$\quad
     BSc Mathematics of Artificial Intelligence}\\[2mm]
    {\color{aimlrule}\rule{\textwidth}{0.8pt}}\\[4mm]
    {\color{aimlmuted}\large Lecture #1 \quad$\cdot$\quad Part V \quad$\cdot$\quad
     Extended notes}\\[3mm]
    {\Huge\bfseries\color{aimlprimary}\baselineskip=1.15\baselineskip #2\par}\vspace{4mm}
    {\color{aimlmuted}\large Dataset: #3}\\[3mm]
    {\color{aimlrule}\rule{\textwidth}{0.8pt}}
  \end{flushleft}
  \vspace{2mm}
  \begin{keybox}[title={Why these notes exist}]
    Lectures 19--22 sit outside G\'eron, the course's primary text. For those
    four lectures \emph{these notes are the primary source}, and they are
    examined on exactly the same terms as the textbook parts. Everything in
    the slide deck for this lecture is here, with the arguments written out at
    length rather than compressed onto a slide; the numbers are the ones the
    accompanying notebook prints, and every one of them is a measurement on
    the corpus named above rather than a figure quoted from a paper.
  \end{keybox}
  \vspace{1mm}
  {\small\tableofcontents}
  \vspace{2mm}
  \noindent{\color{aimlrule}\rule{\textwidth}{0.4pt}}
  \vspace{1mm}
}


\begin{document}
\notestitlefor{4}{I}{Training models}{Titanic, 891 passengers}{Chapter 4}

\section{Why a closed form is not enough}

Lecture 2 derived $\hat{\boldsymbol\theta} =
(\mathbf{X}^{\intercal}\mathbf{X})^{-1}\mathbf{X}^{\intercal}y$, valid exactly
when $\mathbf{X}$ has full column rank. Lecture 3 measured a classifier
properly. Both lectures called \code{.fit()} and moved on. Today we open it.

\begin{center}
\small
\begin{tabular}{p{0.32\textwidth}p{0.56\textwidth}}
\toprule
\textbf{Situation} & \textbf{What the normal equation does} \\
\midrule
Many features    & inverts an $(n{+}1)\times(n{+}1)$ matrix --- cost grows like
                   $n^{3}$ \\
Many instances   & forms $\mathbf{X}^{\intercal}\mathbf{X}$, so every row must
                   be seen at once \\
Rank-deficient design & the inverse does not exist \\
A cost that is not squared error & there is no closed form at all \\
\bottomrule
\end{tabular}
\end{center}

The last row is the one that matters. The classifier we build today minimises a
cost whose stationary equations have no solution in elementary functions, so
the answer has to be \emph{walked to}.

\section{Derivation: gradient descent}

You have a cost $J(\boldsymbol\theta)$ and no formula for its minimiser. The
answer must supply three things: a direction to move in, a step size that does
not overshoot, and a reason to believe the point you stop at is the minimum
rather than merely a place where you stopped.

Throughout: $m$ instances, $n$ features, $\mathbf{X}$ of shape
$m \times (n{+}1)$ with a leading column of ones, $\eta > 0$ the
\term{learning rate}, and
\[ J(\boldsymbol\theta) = \mathrm{MSE}(\boldsymbol\theta)
   = \frac{1}{m}\lVert \mathbf{X}\boldsymbol\theta - y\rVert_2^{\,2}. \]

We derive on squared error because it is the one cost whose answer we already
have in closed form --- which makes it the one cost where the method can be
checked against a formula.

\subsection{Which way is downhill}

Move a small distance in a unit direction $\mathbf{u}$. To first order,
\[ J(\boldsymbol\theta + \epsilon\mathbf{u}) - J(\boldsymbol\theta)
   = \epsilon\,\nabla J(\boldsymbol\theta)\cdot\mathbf{u} + o(\epsilon). \]
By Cauchy--Schwarz, $\lvert\nabla J\cdot\mathbf{u}\rvert \le \lVert\nabla
J\rVert$ for every unit $\mathbf{u}$, with equality only when $\mathbf{u}$ is
parallel to $\nabla J$. So
\[ \mathbf{u}^{\ast} = -\frac{\nabla J(\boldsymbol\theta)}
                             {\lVert\nabla J(\boldsymbol\theta)\rVert}. \]
That is why the gradient and not some other vector: it is the \emph{unique}
direction of steepest decrease, and the argument is three lines long.

\subsection{The gradient, and the update}

\begin{derivbox}[title={One partial derivative, then stack them}]
\[
\begin{aligned}
\frac{\partial J}{\partial\theta_j}
 &= \frac{1}{m}\sum_{i=1}^{m}
    2\bigl(\boldsymbol\theta^{\intercal}\mathbf{x}^{(i)} - y^{(i)}\bigr)
    \frac{\partial}{\partial\theta_j}
    \bigl(\boldsymbol\theta^{\intercal}\mathbf{x}^{(i)}\bigr) \\[2pt]
 &= \frac{2}{m}\sum_{i=1}^{m}
    \bigl(\boldsymbol\theta^{\intercal}\mathbf{x}^{(i)} - y^{(i)}\bigr)
    x_j^{(i)},
\end{aligned}
\]
using $\partial(\boldsymbol\theta^{\intercal}\mathbf{x})/\partial\theta_j =
x_j$, because $\sum_k \theta_k x_k$ contains $\theta_j$ exactly once. Read it
as \emph{error $\times$ feature}, averaged over the instances. Stacking the
partials,
\[ \nabla J(\boldsymbol\theta) = \frac{2}{m}\,\mathbf{X}^{\intercal}
   \bigl(\mathbf{X}\boldsymbol\theta - y\bigr), \]
and the update is
\[ \boldsymbol\theta^{(k+1)} = \boldsymbol\theta^{(k)}
   - \eta\,\nabla J\bigl(\boldsymbol\theta^{(k)}\bigr). \]
\end{derivbox}

One matrix product over the whole training set, at cost $O(mn)$ per evaluation
--- linear in the rows, where the normal equation was cubic in the columns. No
matrix is ever inverted. It is called \term{batch} gradient descent because
every step uses the whole training set.

Note that $\eta$ has units: it converts a gradient into a displacement in
parameter space. That is why its right value depends on how the features are
scaled.

\subsection{Where it stops, and whether that is the minimum}

The iteration stops moving exactly when the step is zero:
\[ \nabla J(\hat{\boldsymbol\theta}) = \mathbf{0}
   \iff \mathbf{X}^{\intercal}\mathbf{X}\hat{\boldsymbol\theta}
      = \mathbf{X}^{\intercal}y. \]

\begin{keybox}[title={Descent does not solve a different problem}]
That is the normal equation of Lecture 2. Descent solves the same system
without ever forming $(\mathbf{X}^{\intercal}\mathbf{X})^{-1}$ --- which is why
it survives the rank-deficient case the closed form cannot even state.
\end{keybox}

The Hessian is constant, $\nabla^2 J = \frac{2}{m}\mathbf{X}^{\intercal}
\mathbf{X}$, and positive semi-definite everywhere. So $J$ is convex: no local
minima, no plateaux other than the minimum itself, no saddle points. Descent
cannot get stuck; it can only be slow.

Hold on to the word \emph{convex}. The moment we leave linear models, in
Lecture 9, every one of those guarantees goes.

\subsection{How large may the step be}

Write $H = \nabla^2 J$ and subtract the fixed point from the update:
\[ \boldsymbol\theta^{(k+1)} - \hat{\boldsymbol\theta}
   = \bigl(\mathbf{I} - \eta H\bigr)
     \bigl(\boldsymbol\theta^{(k)} - \hat{\boldsymbol\theta}\bigr). \]
In the eigenbasis of $H$ each coordinate of the error is multiplied by
$(1 - \eta\lambda_i)$ every step, so every mode contracts if and only if
\[ 0 < \eta < \frac{2}{\lambda_{\max}}. \]
Because $J$ is quadratic this is exact, not a local approximation.

\begin{center}
\small
\begin{tabular}{p{0.26\textwidth}p{0.24\textwidth}p{0.36\textwidth}}
\toprule
\textbf{Learning rate} & \textbf{Factor per step} & \textbf{What you see} \\
\midrule
$\eta \ll 2/\lambda_{\max}$ & just under 1 &
  the cost falls slowly, and is still falling when you stop \\
$\eta \approx 1/\lambda_{\max}$ & small &
  the cost drops fast and flattens \\
$\eta > 2/\lambda_{\max}$ & $\lvert 1-\eta\lambda_i\rvert > 1$ &
  the cost rises, oscillates, and overflows to \code{nan} \\
\bottomrule
\end{tabular}
\end{center}

\begin{failbox}[title={Divergence is silent until it is loud}]
A diverging fit does not raise. It returns weights full of \code{inf} and a
model that predicts one value everywhere. Plot the cost against the step
number; it is the only cheap way to see which of the three rows you are in.
\end{failbox}

\subsection{Why you must scale the features}

The step count is governed by the slowest mode. With the best fixed $\eta$ the
error contracts per step by
\[ \frac{\kappa - 1}{\kappa + 1}, \qquad
   \kappa = \frac{\lambda_{\max}}{\lambda_{\min}}
   = \text{cond}\bigl(\mathbf{X}^{\intercal}\mathbf{X}\bigr). \]
$\kappa = 1$ is a circular bowl and one step reaches the bottom; $\kappa = 100$
is a long thin valley the path zig-zags across; and $\kappa$ large because one
column is in thousands and another in units is a \emph{scaling accident},
entirely avoidable.

Standardising the columns is not cosmetic. It changes $\kappa$, and $\kappa$ is
what sets the number of steps.

\subsection{One instance at a time}

Each batch step costs $O(mn)$, unaffordable for large $m$. Draw an index $i$
uniformly at random and use only that row:
\[ \mathbf{g}^{(i)}(\boldsymbol\theta)
   = 2\bigl(\boldsymbol\theta^{\intercal}\mathbf{x}^{(i)} - y^{(i)}\bigr)
     \mathbf{x}^{(i)},
   \qquad
   \mathbb{E}_i\bigl[\mathbf{g}^{(i)}\bigr] = \nabla J(\boldsymbol\theta). \]
The one-instance gradient is an \emph{unbiased} estimate of the batch gradient,
and that one line is what makes stochastic and mini-batch descent legitimate
rather than merely cheap.

\begin{failbox}[title={Unbiased is not the same as right}]
The estimate's variance does not shrink as $\boldsymbol\theta$ approaches
$\hat{\boldsymbol\theta}$: at the minimum the batch gradient is zero and the
individual gradients are not. So the iterates bounce around the minimum
forever, and the remedy is a learning schedule $\eta_k \to 0$.

Averaging over $b$ rows is still unbiased, by the same one-line argument, and
divides the variance by $b$.
\end{failbox}

\begin{center}
\small
\begin{tabular}{lllp{0.30\textwidth}}
\toprule
\textbf{Variant} & \textbf{Rows/step} & \textbf{Cost} & \textbf{Path} \\
\midrule
Batch      & $m$ & $O(mn)$ & smooth, straight to the minimum \\
Mini-batch & $b$ & $O(bn)$ & slightly noisy, and vectorises on hardware \\
Stochastic & 1   & $O(n)$  & very noisy, never settles without a schedule \\
\bottomrule
\end{tabular}
\end{center}

The bouncing is not only a cost. On a non-convex surface --- every network from
Lecture 9 onwards --- it is what lets the iterate leave a bad local minimum.
Here, on a convex bowl, it is pure nuisance.

\begin{keybox}[title={What the derivation buys you}]
That the gradient is error times feature, so a linear model fits without
inverting anything; that convexity licenses trusting the answer from any
starting point, and that this stops being true in Lecture 9; that a rising cost
curve is a step-size fault rather than a bug; that scaling is part of the
algorithm, not tidiness; and that mini-batches are principled.

And one more: the derivation never used the fact that the cost was squared
error until the partial derivative. Swap the cost, redo that one step, and
everything else stands. That is exactly what we do next.
\end{keybox}

\section{Fitting by descent}

\begin{center}
\small
\begin{tabular}{lll}
\toprule
& \textbf{Normal equation / SVD} & \textbf{Gradient descent} \\
\midrule
Cost in $m$      & linear & linear, per step \\
Cost in $n$      & $n^{2.4}$ to $n^{3}$ & linear \\
Out-of-core      & no --- needs all rows & yes, with mini-batches \\
Scaling required & no & yes \\
Hyperparameters  & none & $\eta$, batch size, schedule, stopping rule \\
Other costs      & squared error only & any differentiable cost \\
\bottomrule
\end{tabular}
\end{center}

Few features and data that fits in memory: use the closed form, it has no knobs
to get wrong. Many features, streaming data, or any cost that is not squared
error: descent.

\subsection{Curvature without leaving linearity}

A straight line cannot fit a curve, but
$\hat y = \theta_0 + \theta_1 x + \theta_2 x^2$ is linear in
$\boldsymbol\theta$ even though it is not linear in $x$. Everything derived
above applies unchanged --- the gradient, the convexity, the convergence
condition --- because all of them were statements about $\boldsymbol\theta$.

\code{PolynomialFeatures(degree=d)} also generates every cross product, which
is how a linear model represents an interaction such as ``third class and
travelling alone''. The column count is $\binom{n+d}{d}$, so ``a bit more
capacity'' is never a small change: degree 6 on four numeric columns gives the
model 227 weights to determine from 712 passengers, about three rows per
weight.

\section{The worked example: a probability, not a label}

A marine safety planning unit wants an evacuation-assistance model: for each
passenger profile, how likely that person is to get off unaided --- and they
need to explain the assignment afterwards.

\begin{center}
\small
\begin{tabular}{lp{0.52\textwidth}}
\toprule
\textbf{What they asked for} & \textbf{What that forces} \\
\midrule
``how likely''    & a \emph{probability}, not a label \\
the number is used to allocate & it must be \emph{calibrated} \\
``explain the assignment'' & weights a human can read \\
\bottomrule
\end{tabular}
\end{center}

\begin{keybox}[title={Why requirement 1 is not a formatting preference}]
``This passenger survives'' is not checkable on its own. ``This passenger
survives with probability 0.31'' is: take everyone the model puts near 0.31 and
count.
\end{keybox}

\subsection{The data, and three kinds of hole}

891 passengers, 12 columns. \code{PassengerId} is an index dressed as a
feature: give it to a model and the model will find something in it.

Three columns have holes, and they get three different answers.
\code{Cabin} is missing for 687 passengers --- 77\%, which is not a column with
gaps but a gap with a column --- so do not impute it: take the deck letter
where there is one and give the rest their own level \code{U}, because the
missingness \emph{is} the information. \code{Age} is missing for 177, imputed
by median inside the pipeline so it is recomputed on every training fold.
\code{Embarked} is missing for 2, imputed by the most frequent port; two rows
out of 891 cannot matter, and you should still be able to say what you did.

Note that \code{Deck = "U"} will largely mean third class, because the cabin
number was recorded for passengers whose tickets carried one. So the column is
partly a proxy for \code{Pclass}. Write that down.

\subsection{The label, and a column that looks like free text}

\begin{center}
\small
\begin{tabular}{lrr}
\toprule
\textbf{Outcome} & \textbf{Passengers} & \textbf{Share} \\
\midrule
Did not survive & 549 & 61.6\% \\
Survived        & 342 & 38.4\% \\
\bottomrule
\end{tabular}
\end{center}

Mildly imbalanced --- nothing like the rare-event problem of Lecture 3.
Accuracy will not be absurd here. It will be something more dangerous than
absurd: plausible.

Every name carries a \emph{title}, and the title carries sex, marital status
and --- for \emph{Master} --- being a boy.

\begin{center}
\small
\begin{tabular}{lrr}
\toprule
\textbf{Title} & \textbf{Passengers} & \textbf{Survived} \\
\midrule
Mrs                        & 126 & 79.4\% \\
Miss                       & 185 & 70.3\% \\
Master --- a boy           &  40 & 57.5\% \\
Rare --- Dr, Rev, Col, \ldots &  23 & 34.8\% \\
Mr                         & 517 & 15.7\% \\
\bottomrule
\end{tabular}
\end{center}

\emph{Master} is the only level that is not a re-encoding of \code{Sex}: those
40 boys survived at nearly four times the rate of the men.

Four engineered columns, each with a reason: \code{Title} (a boy is not a man),
\code{FamilySize} (a group evacuates together, and slowly), \code{IsAlone}
(nobody is waiting for you), \code{Deck} (distance to a boat, and a proxy for
class). Note that \code{FamilySize = SibSp + Parch + 1} is an exact linear
function of two columns we already have. Remember it --- we come back to it
with a rank computation, not with an opinion.

\subsection{Split first, and measure the anchor}

Stratified on the \emph{label} this time, not on a predictor: with 179 test
rows an unstratified draw can shift the base rate by several points, and every
downstream number moves with it. That gives 712 training and 179 test
passengers, at rates 0.3834 and 0.3855.

With $p = 0.3834$, a model that has learned nothing except the base rate scores
$1 - p = 0.617$ accuracy and a log loss of
\[ -\bigl(p\log p + (1-p)\log(1-p)\bigr) = 0.666 \]
--- the entropy of the label, in nats. Every number below has to be compared
against it.

\section{Logistic regression}

\subsection{The model, and what the linear part predicts}

\[ \hat p = \sigma\bigl(\boldsymbol\theta^{\intercal}\mathbf{x}\bigr),
   \qquad \sigma(t) = \frac{1}{1 + e^{-t}}. \]
$\boldsymbol\theta^{\intercal}\mathbf{x}$ is exactly the quantity least squares
fitted; everything new is in $\sigma$ and in what we ask $\boldsymbol\theta$ to
minimise. $\sigma$ is strictly increasing and maps $\mathbb{R}$ onto $(0,1)$,
so it never returns exactly 0 or 1 --- which matters when the loss takes a
logarithm of it.

Inverting it:
\[ t = \log\frac{\hat p}{1-\hat p}
     = \boldsymbol\theta^{\intercal}\mathbf{x}. \]

\begin{keybox}[title={The sentence most often got wrong in the examination}]
The linear model does not predict the probability. It predicts the
\emph{log-odds}, and the coefficients act on the \emph{odds}: add 1 to $x_j$
and the odds are multiplied by $e^{\theta_j}$, wherever you started. The
probability has no such rule.
\end{keybox}

\subsection{The loss, and why this one}

\[ J(\boldsymbol\theta) = -\frac{1}{m}\sum_{i=1}^{m}
   \Bigl[y^{(i)}\log\hat p^{(i)}
        + \bigl(1-y^{(i)}\bigr)\log\bigl(1-\hat p^{(i)}\bigr)\Bigr]. \]

Only the probability given to the true class enters each term, and it is
unbounded above: claim 0.001 for a passenger who survived and that term is
$-\log 0.001 \approx 6.9$.

A scoring rule is \term{proper} if the expected score is optimised by reporting
your true belief. Log loss is proper; so is the Brier score
$\frac{1}{m}\sum(\hat p - y)^2$, which is bounded and easier to read, and we
report it too.

\begin{failbox}[title={Accuracy is not proper}]
Under accuracy a model reporting 0.51 scores exactly the same as one reporting
0.99. So accuracy gives a model no reason at all to tell the truth about its
uncertainty, and optimising it destroys calibration.
\end{failbox}

\subsection{Its gradient, and the absence of a closed form}

Redo the one step of the derivation that used the cost.
$\sigma'(t) = \hat p(1-\hat p)$, and
$\partial\ell/\partial\hat p = (\hat p - y)/\hat p(1-\hat p)$, so the two
factors cancel exactly and
\[ \frac{\partial J}{\partial\theta_j}
   = \frac{1}{m}\sum_i \bigl(\hat p^{(i)} - y^{(i)}\bigr)x_j^{(i)}. \]
Error times feature. The same shape as the gradient of the MSE.

And yet setting it to zero gives $\mathbf{X}^{\intercal}(\sigma(\mathbf{X}
\boldsymbol\theta) - y) = \mathbf{0}$, in which $\boldsymbol\theta$ appears
inside a transcendental function; no rearrangement isolates it. But the Hessian
is $\frac{1}{m}\mathbf{X}^{\intercal}\mathbf{S}\mathbf{X}$ with $\mathbf{S} =
\operatorname{diag}(\hat p^{(i)}(1-\hat p^{(i)})) \succeq 0$, so the surface is
a bowl and descent reaches the global minimum.

This is the first model in the course with no formula for its own answer.
Everything in the first half of today is what makes it fittable anyway.

\subsection{The first numbers}

Fitted with \code{C=1e6}, which turns Scikit-learn's default penalty
\emph{off}, so today's coefficients are the data's rather than a compromise
with a penalty. Turning it back on is the next lecture.

\begin{center}
\small
\begin{tabular}{lrr}
\toprule
\textbf{Model} & \textbf{Log loss} & \textbf{Accuracy} \\
\midrule
Report the base rate to everyone & 0.666 & 61.7\% \\
Logistic regression, 10-fold cross-validated & 0.496 & 81.5\% \\
\bottomrule
\end{tabular}
\end{center}

About a quarter below the anchor on log loss and twenty points of accuracy
above the majority class. And the per-fold spread of the log loss is 0.22 ---
larger than every difference we are about to compare. Carry that number.

\section{Three coefficients that cannot be right}

\begin{center}
\small
\begin{tabular}{lrr}
\toprule
\textbf{Feature} & \textbf{Coefficient} & \textbf{Odds multiplier} \\
\midrule
\code{Deck\_T}      & $-5.678$ & $\times 0.003$ \\
\code{Title\_Master}& $+3.484$ & $\times 32.596$ \\
\code{Sex\_female}  & $+3.308$ & $\times 27.325$ \\
\code{Sex\_male}    & $-2.833$ & $\times 0.059$ \\
\code{Title\_Miss}  & $-2.806$ & $\times 0.060$ \\
\code{Title\_Mrs}   & $-2.020$ & $\times 0.133$ \\
\bottomrule
\end{tabular}
\end{center}

\code{Title\_Miss} at $-2.81$, when 70.3\% of them survived against 38.4\%
overall. \code{Title\_Mrs} at $-2.02$, the highest-surviving group in the data.
And \code{Deck\_T} is the largest weight in the model, when there is exactly
one passenger on deck T.

These are not weak effects or noisy estimates. They have the wrong \emph{sign},
and the model's predictions are nonetheless good. Before reaching for a story
about confounding, ask the cheaper question: what is the shape of the design
matrix?

\subsection{Count the columns, then count the rank}

\begin{center}
\small
\begin{verbatim}
print("columns:", Z_b.shape[1])            # 29
print("rank:   ", np.linalg.matrix_rank(Z_b))  # 23
print("cond:   ", np.linalg.cond(Z_b))         # 2.6e+16
\end{verbatim}
\end{center}

Six columns are linear combinations of the others. Five come from the encoder
--- each one-hot block sums to 1 in every row, and the intercept column is also
1 in every row, so five categorical columns give five exact dependencies. The
sixth is ours: \code{FamilySize = SibSp + Parch + 1}, exactly, for every
passenger, with maximum residual $0.0$. Not ``highly correlated'' --- an
identity. Standardising does not help, because an affine map of a linear
dependence is still a linear dependence.

\begin{failbox}[title={What a rank-deficient design does to the fit}]
If $\mathbf{X}\mathbf{v} = \mathbf{0}$ for some $\mathbf{v} \neq \mathbf{0}$,
then $\mathbf{X}(\boldsymbol\theta + c\mathbf{v}) = \mathbf{X}
\boldsymbol\theta$ for every $c$: along $\mathbf{v}$ the loss is flat.

There is not one best $\boldsymbol\theta$ but a whole affine subspace of them.
The solver returns one point of it and says nothing. Predictions are
unaffected; \emph{every claim about a coefficient is void}.
\end{failbox}

Proved on our own matrix: add 2.5 to \code{Sex\_female} and to
\code{Sex\_male}, subtract 2.5 from the intercept, and the largest change in
any of the 712 logits is $1.8\times10^{-15}$. \code{Sex\_male} can be $-2.83$
or $-0.33$, and both fits make the same prediction for every passenger.

\begin{keybox}[title={The test, stated generally}]
If a coefficient can be moved without changing any prediction, that coefficient
is not a property of the data.
\end{keybox}

\subsection{The repair}

Drop one level per categorical block and remove \code{FamilySize}:

\begin{center}
\small
\begin{tabular}{lrrr}
\toprule
\textbf{Design matrix} & \textbf{Columns} & \textbf{Rank} &
\textbf{Condition number} \\
\midrule
As first written & 29 & 23 & $2.6\times10^{16}$ \\
One level dropped, \code{FamilySize} removed & 23 & 23 & 83.6 \\
\bottomrule
\end{tabular}
\end{center}

Fourteen orders of magnitude. \code{IsAlone} stays: it is an indicator, not a
linear function of the counts.

Each dropped level becomes the \term{reference}, folded into the intercept, and
every remaining coefficient is read \emph{relative to it}. A coefficient quoted
with no stated reference level cannot be read by anybody, including you in a
month --- and this is the most common defect in the regression tables you will
be asked to review.

Use \code{min\_frequency} rather than \code{handle\_unknown="ignore"} here:
with \code{drop="first"} an ignored unknown level is encoded all-zeros, which
is the encoding of the reference level, so our one deck-T passenger would be
scored as though they were on deck A whenever a fold holds them out.

\begin{center}
\small
\begin{tabular}{lr}
\toprule
\textbf{Encoding} & \textbf{Cross-validated log loss} \\
\midrule
All levels, \code{FamilySize} kept --- rank 23 of 29 & 0.496 \\
One level dropped, \code{FamilySize} removed --- full rank & 0.468 \\
\bottomrule
\end{tabular}
\end{center}

\begin{keybox}[title={What that table does and does not license}]
Not a large move, and not the point. The point is that the first row was
produced by a fit that had no unique answer, and nothing said so. The fold
spread of the repaired model is 0.13, and of the rank-deficient one 0.22 ---
both far wider than this gap.

Do not claim the second encoding is \emph{better} on this evidence. Claim it is
\emph{defined}.
\end{keybox}

Now a coefficient can be read aloud. \code{Pclass\_3} $= -1.169$: holding every
other recorded feature fixed, travelling third class multiplies the odds of
survival by $e^{-1.169} = 0.31$, relative to first class. Every clause matters
--- ``holding every other feature fixed'' is doing real work, the multiplier is
on odds rather than probability, and it is always relative to the reference
level. And unique is not the same as stable: \code{Sex} and \code{Title} still
carry much of the same information, so how the effect splits between them moves
with the sample. We only fixed uniqueness.

\section{Capacity, and the two ways a fit fails}

A linear term in \code{Age} cannot say ``children and the elderly were helped
first'', because that is not a monotone statement about age. Add
$\text{Age}^2$: the log-odds become quadratic in age, the boundary becomes a
conic, and the model is still linear in $\boldsymbol\theta$.

\begin{center}
\small
\begin{tabular}{rrrrr}
\toprule
\textbf{Degree} & \textbf{Columns} & \textbf{Training} &
\textbf{Held-out} & \textbf{Held-out accuracy} \\
\midrule
1 &  22 & 0.395 & 0.468 & 81.5\% \\
2 &  32 & 0.381 & 0.467 & 82.2\% \\
3 &  52 & 0.355 & 0.538 & 79.6\% \\
4 &  87 & 0.310 & 1.390 & 77.4\% \\
5 & 143 & 0.286 & 1.922 & 76.0\% \\
6 & 227 & 0.284 & 1.836 & 76.7\% \\
\bottomrule
\end{tabular}
\end{center}

Degrees 1 and 2 differ by 0.001 in held-out log loss, and the fold spread is
over a hundred times that. Those two are a \emph{tie}; reporting degree 2 as
the winner would be reporting noise.

From degree 2 to degree 5 held-out accuracy falls by about six points while
held-out log loss is multiplied by more than four. Accuracy only notices when a
prediction crosses the cut-off; log loss notices a confident prediction
becoming a confidently \emph{wrong} one. The failure is not that the model gets
more passengers wrong --- it is that it becomes certain about the ones it gets
wrong. The requirement was for probabilities, and the metric that matches the
requirement is the metric that saw the damage.

\subsection{A warning worth reading}

\begin{center}
\small
\begin{tabular}{lrrrrrr}
\toprule
\textbf{Degree} & 1 & 2 & 3 & 4 & 5 & 6 \\
\midrule
Converged        & yes & yes & yes & no & no & no \\
Iterations used  & 74 & 256 & 935 & 4000 & 4000 & 4000 \\
\bottomrule
\end{tabular}
\end{center}

The reflex is to raise \code{max\_iter}. It cannot help, and understanding why
is the point.

\begin{failbox}[title={Under separation the minimiser does not exist}]
With 87 columns and 712 rows the two classes eventually become linearly
separable. If some $\boldsymbol\theta$ separates them perfectly, so that
$\boldsymbol\theta^{\intercal}\mathbf{x}^{(i)}$ has the right sign for every
$i$, then $J(c\,\boldsymbol\theta) \to 0$ as $c \to \infty$.

The likelihood has no maximum: scaling the weights up always improves it, and
they run away to infinity. The warning is not a numerical nuisance --- it is
the solver reporting that the estimate you asked for \emph{does not exist}.
\end{failbox}

\section{Calibration, and the cut-off nobody chose}

\begin{definition}[Calibration]
$P(y = 1 \mid \hat p(\mathbf{x}) = q) = q$ for every $q \in (0,1)$.
\end{definition}

Take every passenger the model scored near $q$; the fraction who actually
survived should be about $q$. This is a property of a \emph{set} of
predictions, not of any one prediction --- which is exactly why it is so easy
to ship a model that fails it.

Check it out of fold, with \code{cross\_val\_predict} rather than
\code{predict}, so every probability came from a fit that had never seen that
passenger. A reliability diagram drawn on training predictions always looks
excellent.

Read the diagram by looking at the marker area first: a bin with six passengers
has an observed rate that swings by a sixth on one person, so deviations in the
middle bins are mostly sampling noise while deviations in the crowded end bins
are not. Our expected calibration error, out of fold, is 0.050 --- and it is
roughly calibrated because log loss is a proper scoring rule and we minimised
it directly. The calibration is not a happy accident; it is what the objective
was asking for.

\begin{failbox}[title={A bare label hides a decision nobody made}]
\code{.predict()} is \code{.predict\_proba() >= 0.5}. That cut-off is
cost-optimal only when the two mistakes cost the same. Nothing in the code
mentions 0.5, nothing errors, and the accuracy printed is true --- it just
answers a question the stakeholder did not ask.
\end{failbox}

Measured rather than asserted, over twenty splits, choosing the cut-off on
out-of-fold training predictions against a stated 10 : 1 cost ratio:

\begin{center}
\small
\begin{tabular}{lrr}
\toprule
\textbf{Cut-off} & \textbf{Accuracy} & \textbf{Expected cost} \\
\midrule
0.50 --- the library default & 81.8\% & 169 \\
0.86 --- chosen from the costs, per split & 74.1\% & 59 \\
\bottomrule
\end{tabular}
\end{center}

The default costs 110 escort-units more per evacuation on average, and it is
worse on 20 splits out of 20. \emph{Accuracy goes down when you do the right
thing} --- so if accuracy is the headline, you will reject the better rule.

For a calibrated model the cost-minimising cut-off is
$t^{\ast} = c_{\text{FN}}/(c_{\text{FN}} + c_{\text{FP}})$, which is 0.909 at a
10 : 1 ratio. Our measured optimum sits below it, and that gap is the
calibration error expressed in a unit somebody cares about.

So: fit and report on a proper scoring rule, both cross-validated; choose any
cut-off on out-of-fold predictions against costs or a capacity somebody has
stated; and report accuracy at that cut-off as a \emph{consequence}, never as
the objective.

\section{More than two classes}

Softmax regression gives one parameter vector per class:
\[ s_k(\mathbf{x}) = \boldsymbol\theta_k^{\intercal}\mathbf{x},
   \qquad
   \hat p_k = \frac{\exp(s_k(\mathbf{x}))}{\sum_{j=1}^{K}\exp(s_j(\mathbf{x}))}, \]
so the output is a distribution over the $K$ classes by construction. It is
multi\emph{class}, not multi\emph{label}: four objects in one photograph need
$K$ separate binary classifiers, not this.

Its cost is cross-entropy,
\[ J(\boldsymbol\Theta) = -\frac{1}{m}\sum_{i=1}^{m}\sum_{k=1}^{K}
   y_k^{(i)}\log\hat p_k^{(i)}, \]
whose inner sum keeps exactly one term. At $K = 2$ this is the log loss, term
for term: logistic regression is the two-class case of softmax, not a different
model. And its gradient is
\[ \nabla_{\boldsymbol\theta_k}J = \frac{1}{m}\sum_i
   \bigl(\hat p_k^{(i)} - y_k^{(i)}\bigr)\mathbf{x}^{(i)} \]
--- error times feature, for the third time today. That is not a coincidence;
it is a property of this family of losses paired with these output functions,
and it is why one optimiser fits all of them, including every network in the
second half of this course.

The boundary between classes $j$ and $k$ is where
$(\boldsymbol\theta_k - \boldsymbol\theta_j)^{\intercal}\mathbf{x} = 0$: a
hyperplane, so the decision regions are convex polyhedra. Which is also the
limitation --- a class that is not linearly separable from the union of the
others cannot be carved out, however many iterations you run.

\section{What today's model is still missing}

\begin{center}
\small
\begin{tabular}{p{0.36\textwidth}p{0.28\textwidth}p{0.24\textwidth}}
\toprule
\textbf{Left open} & \textbf{What it means} & \textbf{What fixes it} \\
\midrule
Rank 23 of 29 columns & the minimiser is not unique & a penalty makes it
                                                      unique \\
lbfgs will not converge at degree 4 & the minimiser does not exist & a penalty
                                                      makes it exist \\
Held-out log loss quadruples from degree 2 to 5 & capacity with no counterweight
                                                & a penalty buys most of it
                                                  back \\
The degree was chosen by reading the held-out score & a choice with no
                                                diagnosis behind it &
                                                learning and validation
                                                curves \\
\bottomrule
\end{tabular}
\end{center}

The first two are statements about the optimisation problem, not about the
passengers. One idea repairs all four, and it is the next lecture.

\section{The notebook}

\begin{trybox}[title={What to change}]
Gradient descent is implemented from the derivation in a dozen lines, checked
against the normal equation to eight decimals, and run at three learning rates,
one of which diverges.

Raise \code{ETA\_DIVERGE} and watch where the cost first becomes \code{inf}.
Then compute $2/\lambda_{\max}$ for that design matrix and check the two agree
--- which turns the convergence condition from a claim on a page into a
measurement you made.
\end{trybox}

\section{Exercises}

\begin{enumerate}[leftmargin=1.6em,itemsep=4pt]
  \item State the condition on $\eta$ for batch gradient descent on the MSE to
        converge, and explain why the argument is exact rather than a local
        approximation. \hfill \scope{[5]}
  \item A design matrix has 29 columns and rank 23. A colleague reports that
        one coefficient is $-2.83$ and concludes the feature is protective.
        State what is wrong with the conclusion, and give a two-line
        computation that would demonstrate it. \hfill \scope{[5]}
  \item A logistic fit reports \code{ConvergenceWarning} at 4{,}000 iterations
        on a degree-4 polynomial expansion of 712 rows. Explain why raising
        \code{max\_iter} cannot help. \hfill \scope{[4]}
  \item Held-out accuracy falls by six points between two models while held-out
        log loss quadruples. Explain what has happened to the predictions, and
        which metric you would report given a brief that asks for
        probabilities. \hfill \scope{[4]}
  \item A calibrated model is used with a 10 : 1 cost ratio between false
        negatives and false positives. Give the cost-minimising cut-off and
        explain why the default 0.5 produces a \emph{more accurate} and
        \emph{more expensive} rule. \hfill \scope{[4]}
\end{enumerate}

\section*{Summary}
\addcontentsline{toc}{section}{Summary}

Descent moves along $-\nabla J$, which Cauchy--Schwarz identifies as the unique
steepest direction. On the MSE that gradient is
$\frac{2}{m}\mathbf{X}^{\intercal}(\mathbf{X}\boldsymbol\theta - y)$ and its
zero is the normal equation, so descent solves Lecture 2's system without ever
forming an inverse. The cost is convex, so any stationary point is global; the
step converges if and only if $\eta < 2/\lambda_{\max}$; and the step count is
set by the condition number, which is why scaling is part of the algorithm. One
row gives an unbiased gradient, which is what makes mini-batches legitimate.

Logistic regression predicts the log-odds linearly, minimises a proper scoring
rule, and has no closed form --- but its cost is convex, so the machinery above
fits it. Its fit failed in two distinct ways here: \emph{not unique}, because
six of 29 columns were linear combinations of the others, proved by shifting
two coefficients by 2.5 and moving no logit by more than $10^{-15}$; and
\emph{non-existent} at degree 4, where separation sends the weights to
infinity. Repairing the encoding took the condition number from
$2.6\times10^{16}$ to 83.6.

And a probability is checkable where a label is not. \code{.predict()} applies
a cut-off of 0.5 that nobody chose; choosing it against a stated 10 : 1 cost
ratio was worse on accuracy, better on cost, on 20 splits out of 20.

\end{document}
